\section*{1. Parametry emisyjne dźwięku}
Współczesna postprodukcja wymaga rygorystycznego trzymania się norm głośności dla różnych platform dystrybucyjnych.

\section*{2. Zestawienie standardów LUFS}
\begin{tabularx}{\textwidth}{lXX}
\toprule
\textbf{Platforma / Format} & \textbf{Docelowa głośność (Integrated LUFS)} & \textbf{Max True Peak} \\
\midrule
Kino (DCP 5.1) & -24.0 LUFS ($\pm 1$) & -0.5 dBTP \\
Streaming (VOD) & -27.0 LUFS ($\pm 1$) & -1.0 dBTP \\
YouTube / Web & -14.0 LUFS ($\pm 1$) & -1.0 dBTP \\
\bottomrule
\end{tabularx}

\section*{3. Wnioski technologiczne}
Zachowanie odpowiedniego zapasu dynamiki (headroom) zapobiega degradacji miksu podczas kompresji stratnej w serwisach internetowych.

\vspace{1cm}
\noindent\textbf{Zatwierdził:} Tymon Lewandowski